\documentclass{article}
\usepackage{graphicx}
\usepackage{amsfonts}
\usepackage{amsmath}
\usepackage{amsthm}
\usepackage{tikz}

\theoremstyle{definition}
\newtheorem{theorem}{Theorem}[section]
\newtheorem{corollary}[theorem]{Corollary}
\newtheorem{proposition}[theorem]{Proposition}
\newtheorem{lemma}[theorem]{Lemma}
\newtheorem{definition}[theorem]{Definition}

\theoremstyle{remark}
\newtheorem{remark}[theorem]{Remark}
\newtheorem{example}[theorem]{Example}
\newtheorem{claim}[theorem]{Claim}

\title{Integral and the Fundamental Theorem of Calculus}
\author{Sara}
\date{January 2026}

\begin{document}

\maketitle

\section{Upper Bound}

\begin{definition}[Upper Bound]
    Let $A\subseteq \mathbb{R}$ be a subset of $\mathbb{R}$. If $c$ is an upper bound of $A$, we know that $\forall x\in A,\;x\leq c$.
    Furthermore, the least upper bound, or the supremum of $A$ mean:
    \begin{itemize}
        \item $c$ is an upper bound of $A$
        \item for any upper bound $b$, we have $b\geq c$
    \end{itemize}
    $c$ is call the maximum of $A$ if $c$ is an element of $A$.
\end{definition}

\begin{proposition}[Least Upper Bound Principle]
    Let $A\subseteq \mathbb{R}$ be a subset of $\mathbb{R}$.
    If $A$ is bounded above and is not empty, then $A$ has a least upper bound.
\end{proposition}

\section{Supremum and Infimum on a Function}

\begin{definition}
    $\sup_{x\in I}f(x)$
\end{definition}

\begin{theorem}
    Let $f$ be a function defined on the domain $I\neq \emptyset$. If $f$ is bounded above on $I$, then $f$ has a supremum on $I$.

    Similarly, if $f$ is bounded below on $I$, then $f$ has a infimum on $I$.
\end{theorem}

\section{Integral}

\begin{definition}[Partition]
    A partition of the interval $[a,b]$ is a finite subset $P\subseteq[a,b]$ that satisfies $a,b\in P$.
\end{definition}

\begin{definition}
    Let $f$ be a bounded function on $[a,b]$.
    Let $P=\{x_0,x_1,\dots, x_N\}$ be a partition on $[a,b]$.
    For each $i=1,\dots ,N$, let $m_i$ be the infimum of $f$ at $[x_{i-1},x_i]$, $M_i$ be the supremum of $f$ at $[x_{i-1},x_i]$.
    Write $\Delta x_i=x_i-x_{i-1}$.
    Then the P-lower sum is \[
    L_p=\Sigma_{i=1}^N m_i\cdot\Delta x_i
    \]
    P-upper sum is \[
    U_p=\Sigma_{i=1}^N M_i\cdot\Delta x_i
    \]
\end{definition}

\section{Fundamental Theorem of Calculus}

\end{document}